%==============================================================================
% Sjabloon poster bachproef
%==============================================================================
% Gebaseerd op document class `a0poster' door Gerlinde Kettl en Matthias Weiser
% Aangepast voor gebruik aan HOGENT door Jens Buysse en Bert Van Vreckem

\documentclass[a0,portrait]{hogent-poster}

% Info over de opleiding
\course{Bachelorproef}
\studyprogramme{toegepaste informatica}
\academicyear{2025-2026}
\institution{Hogeschool Gent, Valentin Vaerwyckweg 1, 9000 Gent}

% Info over de bachelorproef
\title{Hoe kan een high-level framework worden opgesteld om CA Gen-applicaties succesvol te migreren naar PL/I code?}
\author{Loïc Vandecandelaere}
\email{loic.vandecandelaere@student.hogent.be}
\supervisor{Thomas Parmentier}
\cosupervisor{Didier Marichal (Arcelor Mittal)}

% Indien ingevuld, wordt deze informatie toegevoegd aan het einde van de
% abstract. Zet in commentaar als je dit niet wilt.
\specialisation{Mainframe}
\keywords{Framework, Migratie, PL/I, CA Gen}
\projectrepo{https://github.com/loicvdc/Bachelorproef-Lo-c-Vandecandelaere}

\begin{document}

\maketitle

\begin{abstract}
Veel organisaties vertrouwen nog op bedrijfskritische applicaties die ontwikkeld zijn met CA Gen. Hoewel deze systemen stabiel functioneren, neemt de beschikbare expertise af en worden onderhoud en verdere ontwikkeling steeds moeilijker. Binnen ArcelorMittal is de Batch Annealing Furnace (BAF)-module de laatste CA Gen-component binnen de lijnplanning die nog naar PL/I moet worden gemigreerd. Eerdere migraties werden succesvol uitgevoerd, maar zonder een vastgelegde en herbruikbare methode.

Dit onderzoek ontwikkelt daarom een high-level migratieframework voor de gecontroleerde migratie van CA~Gen-applicaties naar PL/I. De centrale onderzoeksvraag luidt: Hoe kan een high-level framework worden opgesteld om CA~Gen-applicaties succesvol te migreren naar PL/I-code? Het framework bestaat uit vier onderdelen: een analysefase, vertaalrichtlijnen tussen CA Gen en PL/I, best practices en valkuilen, en een validatiestrategie gebaseerd op de Parallel Run-aanpak.

Via een design-science benadering werd het framework iteratief opgebouwd op basis van documentanalyse, literatuurstudie, reverse engineering en een proof-of-concept (PoC). Hiervoor werd een CA Gen-module gemigreerd naar PL/I.

De validatie via een Parallel Run toonde aan dat de gemigreerde PL/I-module functioneel vergelijkbare output oplevert als het CA~Gen-origineel. Vastgestelde verschillen waren beperkt en terug te voeren op detailfouten in de logica. Dit bevestigt dat het framework een solide basis biedt voor gecontroleerde migraties.

Het onderzoek concludeert dat een framework opgebouwd uit analyse, vertaalschema’s, best practices en validatie effectief richting geeft aan het migratieproces van CA Gen naar PL/I. Het biedt ArcelorMittal een gestructureerde, herhaalbare aanpak die het risico op kennisverlies vermindert en de efficiëntie van toekomstige migraties verhoogt. Daarnaast zijn de analyse-, validatie- en best-practicecomponenten taalagnostisch en inzetbaar voor andere legacy-migraties mits aanpassing van de vertaaltabellen en context.
\end{abstract}

\begin{multicols}{2} % This is how many columns your poster will be broken into, a portrait poster is generally split into 2 columns

\section{Introductie}

Dit onderzoek richt zich op de migratie van CA Gen-gebaseerde legacy-applicaties binnen ArcelorMittal, waar de afnemende beschikbaarheid van expertise rond deze verouderde technologie leidt tot toenemende onderhouds- en moderniseringsrisico’s. Specifiek wordt de Batch Annealing Furnace (BAF)-applicatie onderzocht, de laatste lijnplanningsmodule die nog niet van CA Gen naar PL/I is gemigreerd. Omdat eerdere migraties verspreid over verschillende jaren en ontwikkelaars plaatsvonden zonder een gestandaardiseerde aanpak, heeft dit onderzoek als doel een herbruikbaar high-level migratieframework te ontwikkelen. Dit framework biedt richtlijnen voor de analyse, voorbereiding en omzetting van CA Gen-functionaliteit naar PL/I, ondersteund door best practices en validatiemethoden. De bruikbaarheid ervan wordt geëvalueerd via een proof-of-concept waarbij de BAF-module naar PL/I wordt gemigreerd en de output wordt vergeleken met die van de oorspronkelijke CA Gen-versie, aangevuld met een beoordeling door IT-professionals op vlak van duidelijkheid, volledigheid en praktische toepasbaarheid.


\section{Experimenten}

Het framework werd toegepast op de CA Gen-module \texttt{AB\_APP\_GBP\_MAT\_BAF}.
Tijdens demigratie werden CA Gen-constructen systematisch vertaald naar PL/I, bijvoorbeeld door impliciete database-acties te vervangen door expliciete DB2-cursors en de logica regel voor regel om te zetten volgens de opgestelde vertaaltabel. Daarnaast werden enkele bewuste aanpassingen doorgevoerd, zoals het harmoniseren van de exitstate-codes en het toepassen van bestaande naamgevingsconventies. Dit resulteerde in een werkende PL/I module die functioneel dezelfde output genereerde als de CA Gen-module, met kleine verschillen in waardes bij sommige velden wat wijst op fouten in de logica, wat minimale aanpassingen vereist.

\section{Framework}

Het framework bestaat uit vier grote onderdelen. De analysefase, hier wordt de module in kaart gebracht, views, Action Diagrams, databaseafhankelijkheden en USE-relaties worden geïnventariseerd. Dan volgt een tabel met die CA Genconstructen vertaald naar hun PL/I tegenhanger. Vervolgens een sectie best practices en valkuilen, dit is een lijst met veelgemaakte fouten. Tenslotte een voorstel van een validatiestrategie.

\begin{center}
  \captionsetup{type=figure}
  \includegraphics[width=1.0\linewidth]{schema-framework.png}
  \captionof{figure}{High-level schema van het framework}
\end{center}



\section{Conclusies}

Dit onderzoek toont aan dat een succesvol high-level migratieframework voor CA Gen-applicaties naar PL/I kan worden opgebouwd rond vier kerncomponenten: een analysefase voor het ontsluiten van CA Gen-modellen, vertaalrichtlijnen tussen CA Gen-constructen en PL/I-code, best practices voor het vermijden van veelvoorkomende valkuilen, en een validatiestrategie gebaseerd op Parallel Run. De noodzaak van zo’n framework wordt onderbouwd door de beperkingen van CA Gen, zoals afnemende expertise, vendor lock-in, kennisverlies en de paradigmaverschuiving van modelgestuurd naar procedureel programmeren. De praktische bruikbaarheid van het framework werd aangetoond via een proof-of-concept waarin de BAF-module succesvol naar PL/I werd gemigreerd en functioneel equivalente output produceerde. Het framework biedt ArcelorMittal een gestructureerde en herhaalbare aanpak voor toekomstige migraties, vermindert risico’s op kennisverlies en kan, mits aangepaste vertaaltabellen, ook worden toegepast op migraties naar andere programmeertalen.

\section{Toekomstig onderzoek}

Op basis van dit onderzoek worden volgende vervolgstappen aanbevolen: validatie op meerdere modules, vertaaltabel uitbreiden naar andere talen, semi-automatische tooling met AI, UI- en online-ondersteuning en performance-analyse.

\end{multicols}
\end{document}